\documentclass[12pt,a4paper]{article}

% ─── Packages ────────────────────────────────────────────────────────────────
\usepackage[utf8]{inputenc}
\usepackage[T1]{fontenc}
\usepackage[french]{babel}
\usepackage{geometry}
\geometry{margin=2.5cm}

\usepackage{amsmath, amssymb, amsthm}
\usepackage{algorithm}
\usepackage{algpseudocode}
\usepackage{tikz}
\usetikzlibrary{arrows.meta, positioning, shapes.geometric, decorations.markings}
\usepackage{pgfplots}
\pgfplotsset{compat=1.18}
\usepackage{booktabs}
\usepackage{array}
\usepackage{xcolor}
\usepackage{hyperref}
\usepackage{enumitem}
\usepackage{mdframed}
\usepackage{lmodern}
\usepackage{microtype}
\usepackage{caption}
\usepackage{tcolorbox}
\tcbuselibrary{skins, breakable}

% ─── Colors ──────────────────────────────────────────────────────────────────
\definecolor{urgence}{RGB}{180,30,30}
\definecolor{theorie}{RGB}{30,80,160}
\definecolor{pratique}{RGB}{20,130,80}
\definecolor{attention}{RGB}{200,120,0}
\definecolor{bleuléger}{RGB}{230,240,255}
\definecolor{vertléger}{RGB}{230,248,238}
\definecolor{twist}{RGB}{100,80,170}
\definecolor{twistléger}{RGB}{235,230,255}
\definecolor{grisfonce}{RGB}{80,80,80}
\definecolor{grisclair}{RGB}{245,245,245}

% ─── Custom environments ──────────────────────────────────────────────────────
\newtcolorbox{defbox}[1]{
  colback=bleuléger, colframe=theorie, fonttitle=\bfseries,
  title={#1}, breakable, sharp corners=south
}

\newtcolorbox{metierbox}[1]{
  colback=vertléger, colframe=pratique, fonttitle=\bfseries,
  title={#1}, breakable, sharp corners=south
}

\newtcolorbox{alertbox}[1]{
  colback=yellow!10, colframe=attention, fonttitle=\bfseries,
  title={#1}, breakable, sharp corners=south
}

\newtcolorbox{jurybox}[1]{
  colback=red!5, colframe=urgence, fonttitle=\bfseries,
  title={#1}, breakable, sharp corners=south
}

\newtcolorbox{twistbox}[1]{
  colback=twistléger, colframe=twist, fonttitle=\bfseries\large,
  title={#1}, breakable, sharp corners=south
}

\newtcolorbox{codebox}[1]{
  colback=grisclair, colframe=theorie!50, fonttitle=\bfseries\small,
  title={#1}, breakable, sharp corners=south,
  fontupper=\ttfamily\small
}

% ─── Theorem styles ───────────────────────────────────────────────────────────
\theoremstyle{definition}
\newtheorem{definition}{Définition}[section]
\newtheorem{propriete}{Propriété}[section]
\newtheorem{invariant}{Invariant métier}[section]

\theoremstyle{remark}
\newtheorem{remarque}{Remarque}[section]
\newtheorem{exemple}{Exemple}[section]

% ─── Hyperref setup ───────────────────────────────────────────────────────────
\hypersetup{
  colorlinks=true,
  linkcolor=theorie,
  urlcolor=urgence,
  pdftitle={Twist 10 — Réseau coupé, routage avec dernières données},
  pdfauthor={Système de routage critique},
}

% ─────────────────────────────────────────────────────────────────────────────
\begin{document}

% ═══════════════════════════════════════════════════════════════════
%  COUVERTURE
% ═══════════════════════════════════════════════════════════════════
\begin{titlepage}
  \centering
  \vspace*{2cm}
  {\color{twist}\rule{\linewidth}{2pt}}\\[0.5em]
  {\Huge\bfseries Twist 10}\\[0.4em]
  {\Large\bfseries Le réseau est coupé, mais le centre de dispatch\\[0.2em]
  doit continuer à router avec les dernières données connues}\\[0.4em]
  {\color{twist}\rule{\linewidth}{2pt}}\\[2em]

  % Schéma illustratif : centre déconnecté, véhicules en autonomie
  \begin{tikzpicture}[
    centre/.style={rectangle, rounded corners, draw=theorie, fill=bleuléger,
                   minimum width=2.5cm, minimum height=1.2cm, font=\small\bfseries},
    veh/.style={rectangle, rounded corners, draw=pratique, fill=vertléger,
                minimum width=1.8cm, minimum height=0.8cm, font=\small},
    coupure/.style={-{Latex[length=3mm]}, urgence, thick, dashed},
    data/.style={-{Latex[length=2mm]}, attention, thick},
    scale=1.0
  ]
    \node[centre] (dispatch) at (0,0) {Centre de dispatch};
    \node[veh] (V1) at (-3,-2) {Véhicule 1};
    \node[veh] (V2) at (0,-2) {Véhicule 2};
    \node[veh] (V3) at (3,-2) {Véhicule 3};

    \draw[coupure] (dispatch) -- node[above, sloped, font=\tiny\bfseries\color{urgence}] {COUPURE} (V1);
    \draw[coupure] (dispatch) -- (V2);
    \draw[coupure] (dispatch) -- (V3);

    \draw[data, double] (V1) -- node[below, font=\tiny] {cache local} (V2);
    \draw[data, double] (V2) -- node[below, font=\tiny] {cache local} (V3);
    \draw[data, double] (V3) to[bend left=20] node[right, font=\tiny] {re-synchro} (V1);

    \node[font=\tiny\itshape, color=pratique] at (-1, -3.5)
      {Véhicules continuent avec dernières données connues + échanges pair-à-pair};
  \end{tikzpicture}\\[1.5em]

  {\large Note technique — Résilience, modes dégradés et synchronisation différée}\\[1em]
  {\large\color{gray} Extension du moteur de routage pour la continuité de service}\\[2.5em]

  \begin{tabular}{rl}
    \textbf{Sujet :}        & Sujet 07 — Hackverse 2026 \\
    \textbf{Twist :}        & \#10 — Panne réseau complète \\
    \textbf{Durée typique :} & 10 à 30 minutes (micro-coupure) ou heures (catastrophe) \\
    \textbf{Principe :}     & Ne jamais bloquer une mission à cause d’un réseau absent \\
  \end{tabular}

  \vfill
  {\small\color{gray} Document pédagogique — Présentation jury et équipe technique}
\end{titlepage}

% ═══════════════════════════════════════════════════════════════════
\tableofcontents
\newpage

% ═══════════════════════════════════════════════════════════════════
\section*{Introduction}
\addcontentsline{toc}{section}{Introduction}

Tous les twists précédents supposaient que le système d’aide à la décision
avait accès à des données centralisées actualisées en temps réel
(trafic, météo, fermetures, sens interdits, etc.). Mais cette hypothèse
tombe en panne lorsque \textbf{le réseau de communication est coupé}.

\begin{twistbox}{Twist 10 — Le système doit fonctionner en mode dégradé}
Une panne réseau peut survenir à tout moment : fibre coupée par des travaux,
antenne relais saturée, cyberattaque, ou simple panne du datacenter.
Le centre de dispatch ne reçoit plus de mise à jour des véhicules, et
les véhicules ne reçoivent plus d’instructions centralisées.

Pourtant, les interventions ne s’arrêtent pas. Une ambulance engagée
sur la route doit continuer à naviguer. Un nouveau départ de mission
doit être possible, même sans connexion.

\medskip
Le twist 10 impose de concevoir un système \textbf{résilient} :
\begin{itemize}
  \item Chaque véhicule conserve une copie locale du graphe routier
        et des dernières données connues (cache).
  \item En cas de coupure, il passe en \textbf{mode autonome} :
        il calcule ses propres itinéraires avec ces données figées.
  \item Lorsque la connexion revient, les véhicules et le centre
        \textbf{synchronisent} les changements d’état et les nouvelles
        observations terrain.
\end{itemize}
\end{twistbox}

\newpage

% ═══════════════════════════════════════════════════════════════════
\section{Définition du problème — modes de fonctionnement}

\subsection{Architecture nominale (avec réseau)}

Dans le fonctionnement normal, le centre de dispatch centralise :
\begin{itemize}
  \item les mises à jour de trafic (multiplicateurs $\mu$, fermetures) ;
  \item les prévisions météo ($\mathbf{p}(\theta)$) ;
  \item la position GPS des véhicules (toutes les 30 s) ;
  \item les alertes terrain (conducteurs, capteurs).
\end{itemize}
Le centre calcule les affectations (Twist 9) et les itinéraires, puis
les transmet aux véhicules. Chaque véhicule reçoit des instructions
et suit le chemin. En retour, il envoie sa position et ses observations.

\subsection{Panne réseau — types de coupure}

Deux cas principaux :

\begin{enumerate}[leftmargin=2em, itemsep=4pt]
  \item \textbf{Coupure totale centre ↔ véhicules} : plus aucun message
        ne passe. Le centre perd la connaissance des positions ; les
        véhicules ne reçoivent plus de nouveaux itinéraires ni de mises
        à jour des conditions.
  \item \textbf{Coupure uniquement montante ou descendante} : par exemple,
        les véhicules peuvent envoyer leur position (message court) mais
        ne reçoivent pas de réponse (liaison asymétrique). Le centre
        voit les véhicules mais ne peut pas leur parler.
\end{enumerate}

\subsection{Niveaux de mode dégradé}

\begin{defbox}{Niveaux de résilience}
\begin{itemize}
  \item \textbf{Niveau 0 (nominal)} : connexion OK, synchronisation temps réel.
  \item \textbf{Niveau 1 (coupure < 5 min)} : les véhicules continuent
        avec le dernier itinéraire reçu, sans tentative de recalcul.
  \item \textbf{Niveau 2 (coupure prolongée)} : les véhicules passent
        en mode autonome complet (cache local, recalcul local).
  \item \textbf{Niveau 3 (coupure + changement de mission requis)} :
        le véhicule utilise son cache et son A* embarqué pour produire
        un itinéraire de repli, même sans validation centrale.
\end{itemize}
\end{defbox}

\newpage

% ═══════════════════════════════════════════════════════════════════
\section{Cache embarqué — les dernières données connues}

\subsection{Contenu minimal du cache véhicule}

Chaque véhicule maintient une base locale contenant :

\begin{center}
\begin{tabular}{ll}
  \toprule
  \textbf{Élément} & \textbf{Âge maximal admissible (avant mode dégradé)} \\
  \midrule
  Graphe routier orienté (Twist 7) & 24 h (structure statique) \\
  Multiplicateurs $\mu$ par heure & 2 h (trafic horaire typique) \\
  Fermetures connues (Twist 4) & 15 min \\
  Prévision météo $\mathbf{p}(\theta)$ & 30 min \\
  Score $\kappa(z)$ (Twist 8) & 6 h \\
  Affectation en cours & instantanée \\
  Destination courante & instantanée \\
  \bottomrule
\end{tabular}
\end{center}

Le cache est mis à jour périodiquement (toutes les minutes en fond)
lorsque le réseau est disponible. Un champ \texttt{timestamp}
accompagne chaque donnée.

\subsection{Stratégie de vieillissement des données}

En mode dégradé, le système n’a plus accès aux mises à jour fraîches.
Il doit continuer à router avec des données vieillissantes, mais en
intégrant une \textbf{pénalité de confiance} dans les poids.

\begin{defbox}{Poids avec vieillissement et incertitude de connexion}
Le poids effectif d’un arc à l’instant $t$, alors que la dernière
mise à jour date de $t_0$ (avec $\Delta = t - t_0$), devient :
\[
  w_{\text{mode dégradé}}(e, t) = w(e, t_0) \cdot (1 + \alpha \cdot \Delta)
\]
où $\alpha$ est un coefficient d’augmentation de l’incertitude
(par exemple $\alpha = 0.02$ par minute, soit +2\% de pénalité
par minute d’absence de mise à jour). Au bout de 10 minutes,
le poids est majoré de 20\%.
\end{defbox}

Cette pénalité reflète que les données deviennent moins fiables.
Si la coupure dure trop longtemps, le système finit par privilégier
les axes dont le trafic est structurellement faible (car leur
dégradation relative est moindre).

\subsection{Pseudocode de routage local en mode dégradé}

\begin{algorithm}[H]
\caption{Routage local avec cache (embarqué)}
\begin{algorithmic}[1]
\Require Cache local $C$ (contient graphe $G_0$, multiplicateurs $\mu_0$,
         météo $p_0$, etc.), destination $d$, position $p$, instant $t$
\If{réseau disponible}
    \State Utiliser les données centrales à jour
\Else
    \State $\Delta \gets t - C.\text{timestamp}$
    \If{$\Delta < \text{seuil\_critique}$ (par défaut 5 min)}
        \State Utiliser cache sans pénalité (ou pénalité légère)
    \Else
        \State Appliquer facteur de vieillissement $(1 + \alpha\Delta)$
            sur tous les poids $w(e, t_0)$
    \EndIf
    \State Exécuter A* robuste (Twist 6) avec les données du cache
          et les poids modifiés
\EndIf
\State \Return itinéraire
\end{algorithmic}
\end{algorithm}

\newpage

% ═══════════════════════════════════════════════════════════════════
\section{Synchronisation différée — la reprise du réseau}

\subsection{Événements à synchroniser}

Lorsque la connexion est rétablie, le centre et les véhicules
doivent échanger les événements survenus pendant la coupure :

\begin{itemize}
  \item \textbf{Du véhicule vers le centre} :
    \begin{itemize}
      \item position GPS horodatée (pour reconstruction de trajectoire) ;
      \item nouvelles routes découvertes (arcs fantômes confirmés,
            cf. Twist 8) ;
      \item temps réels de parcours des arcs (pour calibration) ;
      \item incidents constatés (accidents, routes bloquées) ;
      \item changement de mission effectué localement (si décision autonome).
    \end{itemize}
  \item \textbf{Du centre vers le véhicule} :
    \begin{itemize}
      \item mises à jour de trafic accumulées (nouveaux multiplicateurs) ;
      \item fermetures survenues ;
      \item nouvelles affectations (si le centre a continué à tourner) ;
      \item alertes météo fraîches.
    \end{itemize}
\end{itemize}

\subsection{Résolution des conflits}

Deux scénarios de conflit peuvent apparaître :

\begin{enumerate}[leftmargin=2em, itemsep=4pt]
  \item \textbf{Le véhicule a modifié sa mission sans accord central.}
        Le centre compare sa propre estimation (basée sur des données
        périmées) avec le compte-rendu du véhicule. Par défaut, la
        décision terrain prévaut, mais l’incident est journalisé pour
        audit.
  \item \textbf{Le centre a émis un ordre pendant la coupure (non reçu).}
        À la reprise, le centre renvoie l’ordre. Si le véhicule a déjà
        effectué une action incompatible, l’ordre est annulé et une
        explication est générée.
\end{enumerate}

\begin{metierbox}{Règle de réconciliation}
Le système applique une \textbf{chronologie vectorielle} : chaque
entité (centre, véhicule) incrémente un compteur d’événements.
Lors de la synchronisation, on fusionne les événements dans l’ordre
temporel global (basé sur les horodatages des machines, corrigés
par un protocole NTP léger). En cas d’égalité, la version du véhicule
est prioritaire pour les observations terrain ; celle du centre pour
les directives stratégiques.
\end{metierbox}

\subsection{Algorithmes de synchronisation}

\begin{algorithm}[H]
\caption{Synchronisation après coupure (côté centre)}
\begin{algorithmic}[1]
\Require Connexion rétablie avec véhicule $v$
\State $v.\text{sync\_request}()$
\State $v$ envoie : $( \text{positions}[t_1..t_2],\;
       \text{incidents\_observés},\; \text{nouvelles\_routes},\;
       \text{mission\_locale} )$
\State Le centre met à jour sa base avec ces informations
\State Le centre envoie à $v$ les mises à jour globales 
       (trafic, météo, affectations)
\State \textbf{Résolution des conflits} (cf. ci-dessus)
\State $v$ recalcule si nécessaire son itinéraire avec les nouvelles données
\State Journaliser la fin de la coupure et la durée
\end{algorithmic}
\end{algorithm}

\newpage

% ═══════════════════════════════════════════════════════════════════
\section{Scénario : coupure lors d’un incident multiple}

\subsection{Contexte — deux ambulances, trois incidents, réseau coupé}

Reprenons le scénario du Twist 9. À 14h22, trois incidents vitaux
sont signalés. Le centre a affecté : $V_1$ vers Bastos (5,8 min),
$V_2$ vers Bépanda (7,2 min), Mvog-Mbi en attente (sera pris par
$V_1$ après Bastos). À 14h25, alors que $V_1$ est à mi-chemin, une
coupure totale réseau survient.

\subsection{Comportement des véhicules en mode dégradé}

\textbf{Véhicule $V_1$ (vers Bastos)} : Il a reçu l’itinéraire complet
avant la coupure. Il continue à suivre le chemin prévu. Son cache
local contient les données de trafic de 14h22. À 14h27, il arrive à
Bastos, prend en charge le patient. Il doit maintenant décider de
partir vers Mvog-Mbi (comme planifié) ou de faire autre chose. Comme
il est en mode dégradé, il utilise son cache pour calculer le trajet
Bastos → Mvog-Mbi (temps estimé 6,5 min). Il exécute cette instruction.

\textbf{Véhicule $V_2$ (vers Bépanda)} : Il continue également. Il
arrive à Bépanda à 14h29. Son cache lui indique que le centre avait
prévu que $V_1$ prenne Mvog-Mbi. Il reste donc sur place (ou retourne
vers le centre) en attendant la reprise réseau.

\subsection{Retour du réseau à 14h35}

À 14h35, la connexion revient. Le centre reçoit les positions :
$V_1$ à Mvog-Mbi (arrivée 14h33), $V_2$ à Bépanda (arrivée 14h29,
patient pris en charge). Le centre détecte que $V_1$ a suivi le plan
initial. Aucun conflit. Les données de trafic sont mises à jour
(congestion à Bastos entre-temps). $V_1$ reçoit l’ordre de retour
vers l’hôpital. $V_2$ reçoit un nouvel incident (non traité ici).

\begin{metierbox}{Leçon du scénario}
La coupure n’a pas arrêté les interventions : les véhicules ont
poursuivi avec leur dernier plan valide. L’absence de recalcul
pendant 10 minutes n’a pas dégradé significativement la performance
car les conditions n’avaient que peu changé (heure creuse). En cas
de coupure longue, le cache vieillissant pénaliserait les axes
connus pour leur variabilité, incitant les véhicules à prendre
des itinéraires plus sûrs.
\end{metierbox}

\newpage

% ═══════════════════════════════════════════════════════════════════
\section{Invariants métier et cas limites}

\begin{invariant}[Continuité absolue du guidage]
Aucune panne réseau ne doit entraîner l’absence d’itinéraire pour
un véhicule déjà en mission. Le cache local garantit qu’au moins
un chemin (potentiellement sous-optimal) est toujours disponible.
\end{invariant}

\begin{invariant}[Détection de l’absence de réseau]
Le système détecte la coupure en moins de 30 secondes (timeout
sur les messages keep-alive) et bascule automatiquement en mode
dégradé. L’opérateur est informé visuellement.
\end{invariant}

\begin{invariant}[Réconciliation non destructive]
Aucune donnée terrain (nouvelle route, temps réel) collectée pendant
la coupure n’est perdue lors de la synchronisation. Les données
centrales plus récentes peuvent écraser le cache, mais les observations
sont conservées pour l’historique et l’apprentissage.
\end{invariant}

\subsection{Cas limite : coupure très longue (> 2 heures)}

Si la coupure dépasse 2 heures, le cache devient très périmé. Le
système applique alors une stratégie de \textbf{repli maximal} :
\begin{itemize}
  \item Tous les poids sont multipliés par un facteur uniforme de $1 + \alpha \cdot \Delta_{\max}$
        (plafonné à 3,0).
  \item Les arcs fantômes (Twist 8) sont désactivés par défaut
        (car leur probabilité d’existence est trop incertaine).
  \item Le conducteur reçoit une notification : « Données très
        anciennes — prudence accrue. Suivre la route si possible. »
\end{itemize}

\subsection{Cas limite : coupure asymétrique (envoi OK, réception KO)}

Le véhicule continue à envoyer sa position, mais ne reçoit pas
d’instructions. Le centre voit les positions mais ne peut pas
communiquer. Dans ce cas, le centre enregistre les ordres qu’il
aurait voulu envoyer. Dès que la liaison descendante est rétablie,
les ordres sont transmis. Pendant la coupure, le véhicule fonctionne
en mode autonome comme décrit.

\newpage

% ═══════════════════════════════════════════════════════════════════
\section{Implémentation dans le moteur existant}

\subsection{Nouveaux composants}

\begin{center}
\begin{tabular}{>{\ttfamily}p{4cm} p{9cm}}
  \toprule
  \textbf{Fichier} & \textbf{Rôle} \\
  \midrule
  resilience/cache.py &
    Gestion du cache embarqué : sérialisation, horodatage, politique
    d’expulsion (LRU), rechargement périodique. \\[0.4em]
  resilience/network\_monitor.py &
    Surveillance de la connexion (ping, heartbeat, timeout).
    Change l’état du système entre \texttt{NOMINAL}, \texttt{DEGRADE\_LIGHT},
    \texttt{DEGRADE\_FULL}. \\[0.4em]
  resilience/stale\_weight.py &
    Implémente la fonction de pénalité $(1 + \alpha \Delta)$ pour
    les poids en mode dégradé. \\[0.4em]
  resilience/sync.py &
    Protocole de synchronisation après coupure : échange de diffs,
    résolution de conflits, mise à jour des bases. \\[0.4em]
  routing/algorithm.py &
    Ajout d’un paramètre \texttt{use\_cache\_if\_disconnected=True}
    dans \texttt{compute\_route}. \\[0.4em]
  api/views.py &
    Nouvel endpoint \texttt{/api/status/} pour interroger l’état
    du réseau et la fraîcheur du cache. \\
  \bottomrule
\end{tabular}
\end{center}

\subsection{Exemple d’état du cache dans l’interface}

\begin{verbatim}
{
  "network_status": "DEGRADE_FULL",
  "last_contact_center_sec": 187,
  "cache_age": {
    "graph": "2026-04-26T08:00:00Z (age 2h)",
    "traffic": "2026-04-26T09:45:00Z (age 15 min)",
    "weather": "2026-04-26T10:00:00Z (age 0 min - fresh)",
    "closures": "2026-04-26T09:30:00Z (age 30 min)"
  },
  "stale_penalty_factor": 1.30,
  "autonomous_mode": true
}
\end{verbatim}

\newpage

% ═══════════════════════════════════════════════════════════════════
\section{Communication au jury}

\begin{jurybox}{Message clé pour le jury}
\textbf{Le réseau peut tomber, mais pas les services d’urgence.}

\medskip
Notre système n’a pas un seul point de fragilité. Chaque véhicule
embarque une copie locale des données essentielles. En cas de coupure,
il continue à naviguer de façon autonome, en pénalisant prudemment
les informations vieillissantes. Quand la connexion revient, tout
se resynchronise sans perte de données. C’est la \textbf{résilience
par conception}.
\end{jurybox}

\subsection{Questions anticipées}

\paragraph{Question : « Que se passe-t-il si la coupure dure plusieurs heures
et que des routes changent de sens (couloirs de bus horaires) ? »}
\textbf{Réponse :} Les données de sens interdits dynamiques font partie
du cache. Si le cache a plus de 2 heures, le système applique une
pénalité conservative et affiche un avertissement au conducteur.
Dans les faits, les modifications de sens sont rares et connues à
l’avance ; on peut d’ailleurs coder dans le cache une règle temporelle
explicite (ex. « sens unique de 7h à 9h ») qui reste valable même
sans connexion.

\paragraph{Question : « Le cache prend de la place mémoire. Est-ce
réalisable sur un terminal embarqué ? »}
\textbf{Réponse :} Oui. Le graphe routier d’une ville de 2 millions
d’habitants tient en moins de 200 Mo (nœuds et arcs comprimés). Une
tablette ou un smartphone moderne peuvent stocker cela sans difficulté.
Les données de trafic horaire (2\,400 lignes) sont négligeables.

\paragraph{Question : « N’y a-t-il pas un risque de split-brain
(prises de décision contradictoires entre centre et véhicules) ? »}
\textbf{Réponse :} Tout à fait. C’est pourquoi en mode dégradé,
les véhicules n’initient \textbf{jamais} de nouvelle mission sans
ordre central, sauf si cette mission était déjà planifiée (enchaînement
prévu). Pour les décisions critiques (ex. dérivation d’urgence), ils
peuvent le faire mais journalisent l’événement. La réconciliation
post-coupure détecte et résout les conflits par priorité temporelle.

\subsection{Tableau de bord — log de coupure}

\begin{center}
\begin{tabular}{ll}
  \toprule
  \textbf{Champ} & \textbf{Valeur} \\
  \midrule
  Début coupure & 2026-04-26 14:25:03 \\
  Fin coupure   & 2026-04-26 14:35:22 \\
  Durée         & 10 min 19 s \\
  Mode véhicule & dégradé complet (autonome) \\
  Pénalité appliquée & $1 + 0.02 \times 10 = 1.2$ \\
  Itinéraires recalculés localement & 2 (Bastos → Mvog-Mbi) \\
  Conflits détectés & 0 \\
  Données synchronisées : & \\
  - nouvelles routes & 1 (arc fantôme confirmé) \\
  - écarts de temps & 3 arcs mis à jour \\
  - positions GPS & 124 points \\
  \bottomrule
\end{tabular}
\end{center}

\newpage

% ═══════════════════════════════════════════════════════════════════
\section*{Conclusion}
\addcontentsline{toc}{section}{Conclusion}

Le twist 10 est le dernier rempart contre la défaillance totale
du système. Il rappelle une vérité fondamentale : \textbf{un système
critique ne peut pas dépendre d’une infrastructure réseau fiable
à 100\%}. Il doit être conçu pour \emph{continuer à servir} même
quand le monde extérieur se déconnecte.

\medskip
Trois contributions majeures :

\begin{enumerate}[leftmargin=2em, itemsep=4pt]
  \item \textbf{Cache embarqué intelligent} : le véhicule garde une
        copie des dernières données valides et les utilise de façon
        transparente.
  \item \textbf{Pénalisation du vieillissement} : plus les données
        sont âgées, moins on leur fait confiance — sans jamais les
        jeter brutalement.
  \item \textbf{Synchronisation différée} : le système rattrape son
        retard à la moindre opportunité de connexion, sans perte
        d’observation terrain.
\end{enumerate}

\medskip
Ainsi, du twist 1 (graphe statique) au twist 10 (déconnexion), le
système a appris à gérer l’incertitude sous toutes ses formes :
trafic, météo, équité cartographique, topologie, rareté des véhicules,
et enfin panne du réseau. Chaque couche ajoute une brique de robustesse
pour que le véhicule d’urgence arrive toujours — même dans le pire
des mondes.

\medskip
\begin{center}
  {\color{twist}\rule{0.5\linewidth}{1pt}}\\[0.5em]
  {\itshape « Le réseau n’est pas une option, c’est une commodité.
  Un système d’urgence se doit de fonctionner même sans lui. »}\\[0.5em]
  {\color{twist}\rule{0.5\linewidth}{1pt}}
\end{center}

\end{document}